\documentclass[11pt,a4paper]{article}
\usepackage[margin=2cm]{geometry}
\usepackage{amsmath}
\usepackage{siunitx}
\usepackage{booktabs}
\usepackage{graphicx}
\usepackage{float}
\usepackage{pgfplots}
\usepackage{pgfplotstable}
\pgfplotsset{compat=newest}
\sisetup{separate-uncertainty=true}

% ============================================================
%  EDIT PER PRACTICAL: title, name, axis labels, fit endpoints
% ============================================================
\title{Practical: TITLE HERE}
\author{Ibrahim}
\date{\today}

% Axis labels and ranges
\newcommand{\xlabeltext}{$L$ / m}
\newcommand{\ylabeltext}{$T^2$ / s$^2$}
\def\xmin{0}   \def\xmax{1.1}
\def\ymin{0}   \def\ymax{4.5}

% Worst-fit lines: read two points for each line off your graph.
% Max gradient: through the LOWER end of the first error bar and
% the UPPER end of the last one. Min gradient: the opposite.
\def\maxxa{0.20} \def\maxya{0.77} \def\maxxb{1.00} \def\maxyb{4.10}
\def\minxa{0.20} \def\minya{0.87} \def\minxb{1.00} \def\minyb{4.00}

% Gradients calculated from the endpoints above
\pgfmathsetmacro{\gmax}{(\maxyb-\maxya)/(\maxxb-\maxxa)}
\pgfmathsetmacro{\gmin}{(\minyb-\minya)/(\minxb-\minxa)}
\pgfmathsetmacro{\gmean}{(\gmax+\gmin)/2}
\pgfmathsetmacro{\gunc}{(\gmax-\gmin)/2}

\begin{document}
\maketitle

\section{Aim}
State what you are investigating and what you expect to find.

\section{Variables}
\begin{description}
  \item[Independent:] 
  \item[Dependent:] 
  \item[Control:] 
\end{description}

\section{Method}
Describe the apparatus and procedure so someone else could repeat it.

\subsection{Risk assessment}
\begin{table}[H]\centering
\begin{tabularx}{\textwidth}{@{}l X X@{}}
\toprule
Risk & Hazard & Mitigation(s) \\
\midrule
Burns & Hot water bath & Wear goggles; use tongs; keep the bath away from the edge \\
Cuts & Broken glass tube & Check the tube before use; dispose of broken glass safely \\
\bottomrule
\end{tabularx}
\end{table}

\section{Results}
\begin{table}[H]\centering
\pgfplotstabletypeset[
  col sep=comma,
  columns={L,Lerr,T2,T2err},
  columns/L/.style={column name={$L$ / \unit{\metre}}, fixed, fixed zerofill, precision=2},
  columns/Lerr/.style={column name={$\Delta L$ / \unit{\metre}}, fixed, fixed zerofill, precision=3},
  columns/T2/.style={column name={$T^2$ / \unit{\second\squared}}, fixed, fixed zerofill, precision=2},
  columns/T2err/.style={column name={$\Delta T^2$ / \unit{\second\squared}}, fixed, fixed zerofill, precision=2},
  every head row/.style={before row=\toprule, after row=\midrule},
  every last row/.style={after row=\bottomrule}
]{data.csv}
\caption{Results table.}
\end{table}

\begin{figure}[H]\centering
\begin{tikzpicture}
\begin{axis}[
  width=0.85\textwidth, height=8cm,
  xlabel={\xlabeltext}, ylabel={\ylabeltext},
  xmin=\xmin, xmax=\xmax, ymin=\ymin, ymax=\ymax,
  grid=both, minor tick num=1,
  legend pos=north west,
]
\addplot[only marks, mark=x, error bars/.cd, y dir=both, y explicit]
  table[col sep=comma, x=L, y=T2, y error=T2err]{data.csv};
\addlegendentry{Data}

\addplot[blue, no marks] table[col sep=comma, x=L,
  y={create col/linear regression={y=T2}}]{data.csv};
\addlegendentry{Best fit}

\addplot[red, dashed, no marks] coordinates {(\maxxa,\maxya) (\maxxb,\maxyb)};
\addlegendentry{Max gradient}
\addplot[green!50!black, dashed, no marks] coordinates {(\minxa,\minya) (\minxb,\minyb)};
\addlegendentry{Min gradient}
\end{axis}
\end{tikzpicture}
\caption{Graph of $T^2$ against $L$ with error bars, best-fit line and worst-fit lines.}
\end{figure}

\section{Analysis}
Gradient of maximum line: $\num[round-mode=places,round-precision=2]{\gmax}$\\
Gradient of minimum line: $\num[round-mode=places,round-precision=2]{\gmin}$\\
Mean gradient: $\num[round-mode=places,round-precision=2]{\gmean}$\\
Uncertainty in gradient: $\pm\,\num[round-mode=places,round-precision=2]{\gunc}$

Show the equation you are using, rearrange it, and calculate your final result:
\begin{equation}
  T^2 = \frac{4\pi^2}{g}\,L \quad\Rightarrow\quad g = \frac{4\pi^2}{\text{gradient}}
\end{equation}

\section{Uncertainties}
Percentage uncertainties, propagation, and final answer with its uncertainty.

\section{Evaluation}
Sources of error, improvements, and how confident you are in the result.

\end{document}
